% Part: sets-functions-relations
% Chapter: arithmetization
% Section: reflections

\documentclass[../../../include/open-logic-section]{subfiles}

\begin{document}

\olfileid[tr]{sfr}{arith}{ref}
\olsection{Bazı Felsefi Düşünceler}

Teknik ayrıntılar bu kadar. Peki bunlar neyi başardı?

Neredeyse tartışmasız biçimde, bize çok {güzel} bir saf matematik sundular.
Üstelik bazı derin kavramsal başarılar da elde edildi. Tamlık Özelliği'nin gerçel
sayılarla rasyonel sayılar arasındaki belirleyici farkı ifade ettiğini görmek
derin bir kavrayıştı. Dahası, gerçel sayıların Dedekind kesimleri olarak açıkça
kurulması, analizin konusunu sağlam bir temele oturtur. \emph{Tam sıralı
cisim} kavramının tutarlı olduğunu biliyoruz; çünkü kesimler tam da böyle bir
cisim oluşturur.

Bütün bunlara karşın, bu başarılarla ilgili birkaç çekincemizi dile
getirmeliyiz.

İlk olarak gerçel sayıları kesimler aracılığıyla düşünmenin, onları alışıldık
(muhtemelen sonsuz) ondalık açılımları aracılığıyla düşünmekten
\emph{daha} titiz olduğu açık değildir. Gerçel sayıların bu ikinci “kuruluşu”,
onların Cauchy dizileri yoluyla kuruluşuna bazı yönlerden benzer; fakat aslında
17. yüzyılın başlarından itibaren matematikçilerce özünde biliniyordu
(bkz.\ \olref[cauchy]{sec}). Titizlikteki asıl artış, gerçel sayıların Tamlık
Özelliği'ne sahip olduğunun fark edilmesiyle geldi; gerçel sayıları belirli
kümeler olarak kurabilme olanağı kendi başına o kadar da ilginç olmayabilir.

Tam sayıların ya da rasyonel sayıların (çok daha kolay olan)
aritmetikleştirilmesinin bu alanlarda titizliği artırdığı ise daha da
belirsizdir. Burada basit bir gözlem yapmak yararlı olur. Tam sayıları doğal
sayıların sıralı ikililerinin denklik sınıfları olarak \emph{kurduktan},
rasyonel sayıları tam sayıların sıralı ikililerinin denklik sınıfları olarak
kurduktan ve gerçel sayıları rasyonel sayı kümeleri olarak kurduktan hemen sonra
bu kuruluşları \emph{unuturuz}. Özellikle, kuruluşların gerektiği gibi
davrandığını gösteren ispatlar dışında hiç kimse bir matematiksel ispatta bu
kuruluşlara \emph{başvurmak} istemez. Bir gerçel sayıdan doğrudan söz etmek,
doğal sayıların kümelerinin kümelerinin kümelerinin kümelerinin kümelerinin
kümelerinden oluşan bir kümeden söz etmekten çok daha kolaydır.
%(For reals were constructed as sets of rationals; and these were constructed as equivalence classes of ordered pairs of integers, i.e., sets of sets of sets of integers; etc.).
%2/3 = [2,3]
%	i.e. {<2,3>, ... }
%	{ {{2},{2,3}}, ... }
%
%reals are 
%	sets of rationals 
%
%rationals are 
%	sets of sets of sets of integers
%
%integers are
%	sets of sets of sets of naturals
	
Bu tanımların tam sayıların, rasyonel sayıların ya da gerçel sayıların
\emph{metafizik bakımdan} ne olduğunu söylediği ise en kuşkulu noktadır. Başka
bir deyişle gerçel sayıların (örneğin) belirli kümeler (kümelerin kümelerinin
kümeleri\ldots) \emph{olduğu} kuşkuludur. Böyle bir görüşün önündeki temel
engel, kuruluşun birçok farklı biçimde yapılabilmesidir. Gerçel sayılar
durumunda gerçekten ilginç derecede farklı bazı kuruluşlar vardır
(bkz.\ \olref[cauchy]{sec}). Ancak farklı kuruluşlar elde etmenin bütünüyle
önemsiz bir yolu da şudur: \olref[sfr][rel][ref]{sec}'te olduğu gibi sıralı
ikilileri biraz farklı tanımlayabilirdik; bu alternatif sıralı ikili kavramını
kullansaydık kuruluşlarımız yine aynı ölçüde iyi çalışacak, fakat elimizde
farklı nesneler olacaktı. Bu nedenle tam sayıların, rasyonel sayıların ve gerçel
sayıların birbiriyle yarışan birçok küme-teorik kuruluşu vardır. Şimdi tam
sayıların (örneğin) \emph{şu} kümeler değil de \emph{bu} kümeler olduğunu
ileri sürmek yalnızca keyfî (ve mahcup edici) olurdu.
(\olref[sfr][rel][ref]{sec}'te olduğu gibi bu, \citealt{Benacerraf1965}'ın
ünlü kıldığı bir argüman örneğidir.)

Bir nokta daha anılmaya değer: kuruluşlarımızda oldukça \emph{tuhaf} bir şey
vardır. Doğal sayılarla başladık. Sonra tam sayıları ve “tam sayıların
$0$'ını”, yani $\equivrep{0,0}{\Intequiv}$'yi kurduk. Fakat
$0 \neq \equivrep{0,0}{\Intequiv}$'dir. Hatta kuruluşlarımıza göre
\emph{hiçbir} doğal sayı bir tam sayı değildir. Oysa bu, sezgiye son derece
aykırı görünür. Dahası, \olref[sfr][set][imp]{sec}'te fazla gerekçe vermeden
$\Nat \subseteq \Rat$ olduğunu ileri sürdük. Kuruluşlar bize sayıların tam
olarak \emph{ne olduğunu} söylüyorsa bu iddia apaçık yanlıştı.

Öyleyse geri çekilip baktığımızda nereye varıyoruz? Naif bir küme teorisi
içinde çalışıp doğal sayıları hazır kabul ederek tam sayıları, rasyonel
sayıları ve gerçel sayıları belirli kümeler \emph{olarak ele alabiliriz}. Bu
anlamda bu varlıkların teorilerini bir küme teorisi içine \emph{gömebiliriz}.
Ancak bu gömmenin felsefi önemi hiç de basit değildir.

Elbette bunların hiçbiri son söz değildir!{} Vurgulanmak istenen yalnızca
şudur: gerçel sayıların aritmetikleştirilmesinin derin bir felsefi öneme
\emph{sahip} olduğunu göstermek, ek bir \emph{felsefi} argüman gerektirir.

%Additionally: for the entire duration of this chapter, we have treated
%the natural numbers as simply \emph{given} to us. But what can we do
%with them? That is the subject matter for the next chapter.

\end{document}

